% ============================================================
% gravity_as_clock_density.tex
% STATUS: STUB
% Corresponds to: src/experiments/exp_02_gravity_clock_density.py
% ============================================================

% Topics to cover:
% - Information received locally determines clock count
%   - Sparse information (distant star): few clocks, baseline tick rate
%   - Dense information (near mass): many clocks, high scheduler load
%   - Information overload (black hole): processor saturation, time halts
% - Clock density gradient = gravitational potential gradient
%   - No curved geometry required
%   - Refractive bias emerges from differential tick rates across wavefront
% - Gravitational time dilation as scheduler load
%   - Derivation: how clock count maps to observed time dilation
%   - Must reproduce Schwarzschild metric in the continuum limit
% - The event horizon as computational deadlock
%   - Queue depth -> infinity
%   - Hawking radiation as stochastic queue drop (Gambler's Ruin)
%   - Black hole information paradox: information is queued, not destroyed
% - Connection to holographic principle \cite{maldacena1998}
% - Bekenstein-Hawking entropy \cite{bekenstein1973} \cite{hawking1975}
%   as boundary clock count at saturation
% - Verlinde entropic gravity \cite{verlinde2011}
%   -- comparison and distinction
% - Dark matter as clock density anomaly at galactic scales

\textbf{[STUB]}
